\subsection{泰勒公式}

2.5.1. 设 \(r\) 是满足 r \(\geqslant 1\) 的整数，\(f\) 是从 E 的开集 U 到 F 的 \(C^r\) 类映射。对于 \(x \in U\)、\(h \in E\) 和 \(p \leqslant r\)，约定用 \(D^{p}f(x_0).h^p\) 代替 \(D^{p}f(x_0).(h,\ldots,h)\)。若线段 \([x,x + h]\) 包含于 U，则有公式（“泰勒公式”）：

\[
f (x + h) = \sum_ {p = 0} ^ {r - 1} \frac {1}{p !} \mathbf {D} ^ {p} f (x). h ^ {p} + v _ {r} (x; h)
\]

其中“余项” \(v_{r}(x; h)\) 由下式给出：

\[
v _ {r} (x; h) = \int_ {0} ^ {1} \frac {(1 - t) ^ {r - 1}}{(r - 1) !} \mathrm{D} ^ {r} f (x + t h). h ^ {r} d t
\]

有：

\[
v _ {r} (x; h) \equiv \frac {1}{r !} \mathrm{D} ^ {r} f (x). h ^ {r} \quad \text {   mod   } o (\| h \| ^ {r}) \quad \text {   when   } h \text {   tends   to   } 0
\]

以及

\[
f (x + h) \equiv \sum_ {p = 0} ^ {r} \frac {1}{p !} \mathrm{D} ^ {p} f (x). h ^ {p} \quad \mod o (\| h \| ^ {r})
\]

当 h 趋于零时。

2.5.2. 再设 \(\gamma\) 是 \(F\) 上的连续半范数；若对线段 \(\|D^{r}f(z)\|_{\gamma}\leqslant M\) 的每一点 \(z\) 都有 \([x,x + h]\)，则有：

\[
\| v _ {r} (x; h) \| _ {\gamma} \leqslant \frac {\mathbf {M}}{r !} \| h \| ^ {r}
\]

2.5.3. 再假设 \(\mathbf{E} = \mathbf{R}^n\)。于是有：

\[
f (x + h) \equiv \sum_ {| \alpha | \leqslant r} \Delta^ {\alpha} f (x) h ^ {\alpha} \quad \mod o (\| h \| ^ {r}) \text {   when   } h \text {   tends   to   } 0
\]

其中置：

\[
\Delta^ {\alpha} f (x) = \frac {1}{\alpha !} \partial^ {\alpha} f (x)
\]

2.5.4. 设 \(f\) 和 \(g\)\(C^r\) 是 E 的开集 \(U\)\(E\) 上取值于 \(F\) 的两个 C＾r 类函数。\(f\) 和 \(g\) 在 U 中一点 \(x\)\(U\) 处的接触阶 \(\geqslant r\) 的充要条件是，对每个满足 \(D^p f(x) = D^p g(x)\) 的整数 \(p\)，都有 \(0 \leqslant p \leqslant r\)。当 \(E\)\(\leqslant r\) 有限维时，这就是说 \(f\) 和 \(g\) 的  r 阶迭代偏导数（相对于 \(E\) 的一组基）在点 \(x\) 处相等。

2.5.5. 设 U 是 \(E \times R^{n}\) 的一个开子集，形如 \(V \times I_{1} \times \cdots \times I_{n}\)，其中 V 是 E 的开子集，\(I_{1}, \ldots, I_{n}\) 是包含 0 的 R 的开区间。置 \(U_{0} = V\)，并对 \(U_{j} = V \times I_{1} \times \cdots \times I_{j}\) 置 \(1 \leqslant j \leqslant n\)。给定函数 \(f \in \mathcal{C}^{r}(U; F)\)（\(1 \leqslant r \leqslant \infty\)），存在唯一的函数序列 \(f_{j}\in\mathcal{C}^{r-1}(\mathbf{U}_{j};\mathbf{F})\)（\(0\leqslant j\leqslant n\)），使得：

\[
f (x, t _ {1}, \dots , t _ {n}) = f _ {0} (x) + \sum_ {j = 1} ^ {n} t _ {j} f _ {j} (x, t _ {1}, \dots , t _ {j})
\]

对 \(x \in V\) 和 \(t_j \in I_j\) 成立。有：

\[
f _ {0} (x) = f (x, 0, \dots , 0)
\]

\[
f _ {j} (x, t _ {1}, \dots , t _ {j}) = \int_ {0} ^ {1} \partial_ {j} f (x, t _ {1}, \dots , t _ {j - 1}, t _ {j} u, 0, \dots , 0) d u
\]

对 \(1 \leqslant j \leqslant n\) 成立。在最后一个公式中，\(\partial_j f\) 表示函数 (t＿1, , t＿n)  f(x, t＿1, , t＿n) 的第 \(j\)\((t_1, \ldots, t_n) \mapsto f(x, t_1, \ldots, t_n)\) 个偏导数。

